\documentclass{article}
\usepackage[english]{babel}
\usepackage{geometry,amsmath,amssymb,amsthm}
\geometry{letterpaper, margin=1.2in}

\newtheorem{theorem}{Theorem}
\newtheorem{lemma}[theorem]{Lemma}
\newtheorem{remark}[theorem]{Remark}

\newcommand{\C}{\mathbb{C}}
\newcommand{\R}{\mathbb{R}}
\newcommand{\Real}{\operatorname{Re}}
\newcommand{\Imag}{\operatorname{Im}}
\newcommand{\Res}{\operatorname*{Res}}

\begin{document}

\title{Exact Closed-Form European Call Price under Rough Heston\\[4pt]
\large via Three-Term Orthogonal Pad\'e of the Cumulant Generator}
\author{Stephen Crowley}
\date{May 2026}
\maketitle

\tableofcontents

% ------------------------------------------------------------------
\section{Model and notation}
% ------------------------------------------------------------------

Let $S_t$ follow the rough Heston model with Hurst exponent
$H\in(0,\tfrac{1}{2})$, $\mu:=H+\tfrac{1}{2}$,
mean-reversion speed $\lambda>0$, long-run variance $\theta>0$,
vol-of-vol $\nu>0$, correlation $\rho\in(-1,1)$, initial variance
$V_0>0$, and risk-free rate $r\ge0$. The Riccati coefficients are
\begin{equation}
P(v)=\tfrac{1}{2}(-v^2-iv),\qquad
Q(v)=\lambda(iv\rho\nu-1),\qquad
R=\frac{\lambda^2\nu^2}{2},
\end{equation}
and the log-moneyness is
\begin{equation}
k=\log(K/S_0)-rT.
\end{equation}
The Riemann--Liouville fractional integral and derivative of order
$\mu\in(\tfrac{1}{2},1)$ are
\begin{align}
I^r f(t) &:= \frac{1}{\Gamma(r)}\int_0^t(t-s)^{r-1}f(s)\,ds,\\
D^\mu f &:= \frac{d}{dt}\bigl(I^{1-\mu}f\bigr).
\end{align}
The log-characteristic function of $X_T:=\log(S_T/S_0)-rT$ is
\begin{equation}
\mathbb{E}\bigl[e^{ivX_T}\bigr]=\exp\bigl(\Phi(v,T)\bigr),
\end{equation}
where
\begin{equation}\label{eq:Phi_def}
\Phi(v,T)=\lambda\theta\int_0^T h(v,s)\,ds
+V_0\,I^{1-\mu}h(v,\cdot)(T),
\end{equation}
and $h(v,\cdot)$ is the unique solution of the fractional Riccati
initial-value problem
\begin{equation}\label{eq:Riccati}
D^\mu h = P(v)+Q(v)h+R\,h^2,\qquad I^{1-\mu}h\big|_{t=0}=0.
\end{equation}

% ------------------------------------------------------------------
\section{M\"untz--Tau coefficients and cumulant moments}
% ------------------------------------------------------------------

\subsection{Recurrence for $h$}

The unique solution of~\eqref{eq:Riccati} admits the M\"untz expansion
\begin{equation}
h(v,t)=\sum_{j=1}^{\infty}a_j(v)\,t^{j\mu},
\end{equation}
where the coefficients $a_j(v)\in\C[v]$ satisfy
\begin{align}
a_1(v) &= \frac{P(v)}{\Gamma(\mu+1)},
\label{eq:a1}\\
a_j(v) &= \frac{\Gamma((j-1)\mu+1)}{\Gamma(j\mu+1)}
\left(Q(v)\,a_{j-1}(v)
+R\sum_{\ell=1}^{j-2}a_\ell(v)\,a_{j-1-\ell}(v)\right),
\quad j\ge2.
\label{eq:aj}
\end{align}
Each $a_j(v)$ has degree $j+1$ and satisfies
\begin{equation}
a_j(0)=a_j(-i)=0.
\end{equation}

\subsection{Cumulant lattice coefficients}

The scalar integration weights are
\begin{equation}
u_k=\begin{cases}0 & k=0,\\[2pt]
\dfrac{\lambda\theta}{k\mu+1} & k\ge1,
\end{cases}
\qquad
w_k=\begin{cases}V_0\,\Gamma(\mu+1) & k=0,\\[4pt]
V_0\,\dfrac{\Gamma((k+1)\mu+1)}{\Gamma(k\mu+2)} & k\ge1.
\end{cases}
\end{equation}
Setting $a_0:=0$, the cumulant lattice coefficients are
\begin{equation}\label{eq:dk}
d_k(v):=u_k\,a_k(v)+w_k\,a_{k+1}(v)\in\C[v],\qquad k\ge0.
\end{equation}
The power rule
\begin{equation}
I^r t^{j\mu}=\frac{\Gamma(j\mu+1)}{\Gamma(j\mu+r+1)}\,t^{j\mu+r},
\end{equation}
applied termwise to both integrals in~\eqref{eq:Phi_def} and
reindexing $j\mapsto k+1$ on the fractional-integral sum, gives
\begin{equation}\label{eq:Phi_series}
\Phi(v,T)=\sum_{k=0}^{\infty}d_k(v)\,T^{k\mu+1}
=:T\cdot\widetilde\Phi(T^\mu;v).
\end{equation}

% ------------------------------------------------------------------
\section{Three-term orthogonal Pad\'e of $\widetilde\Phi$}
% ------------------------------------------------------------------

\subsection{Moment functional and FOPS}

Treat $\widetilde\Phi(z;v)=\sum_{k\ge0}d_k(v)z^k$ as a formal power
series with moment functional
\begin{equation}
\mathcal{L}_v[z^k]:=d_k(v).
\end{equation}
The Chebyshev/Stieltjes algorithm run for $M+1$ steps on $\{d_k(v)\}$
produces TTRR coefficients
\begin{equation}
\alpha_n(v)=\frac{\mathcal{L}_v[z\,p_n^2]}{\mathcal{L}_v[p_n^2]},
\qquad
\beta_n(v)=\frac{\mathcal{L}_v[p_n^2]}{\mathcal{L}_v[p_{n-1}^2]},
\end{equation}
and the monic formal orthogonal polynomial sequence
$\{p_n(z;v)\}_{n=0}^{M+1}$ through
\begin{equation}\label{eq:TTRR}
p_{n+1}=(z-\alpha_n)\,p_n-\beta_n\,p_{n-1},
\qquad p_{-1}\equiv0,\quad p_0\equiv1.
\end{equation}
The associated polynomials of the first kind $\{p_n^{(1)}\}$ satisfy
the same recurrence with initial data
\begin{equation}
p_{-1}^{(1)}\equiv-1,\qquad p_0^{(1)}\equiv0.
\end{equation}

\subsection{The $[M/M]$ Pad\'e element}

By the Pad\'e--FOPS correspondence (Brezinski, Gautschi, Chihara),
\begin{equation}\label{eq:SM}
S_M(z;v):=d_0(v)\cdot\frac{p_{M+1}^{(1)}(z;v)}{p_{M+1}(z;v)}
\end{equation}
is the $[M/M]$ Pad\'e element of $\widetilde\Phi(\cdot;v)$, matching
$2M+2$ moments. Stahl's theorem gives convergence in logarithmic
capacity outside the minimal-capacity singular set, with a-posteriori
bound
\begin{equation}\label{eq:aposteriori}
|\widetilde\Phi-S_M|\le\frac{|\Delta_M|^2}{|\Delta_{M-1}|-|\Delta_M|},
\qquad\Delta_k:=S_k-S_{k-1}.
\end{equation}

% ------------------------------------------------------------------
\section{The $v$-plane partial fraction}
% ------------------------------------------------------------------

Set
\begin{equation}\label{eq:PhiM}
\Phi_M(v,T):=T\cdot S_M(T^\mu;v).
\end{equation}
Since $\alpha_n,\beta_n\in\C(v)$, the function $\Phi_M(v,T)$ is
rational in $v$ and admits the partial fraction
\begin{equation}\label{eq:PhiM_pf}
\Phi_M(v,T)=\sum_{j=1}^{N}\frac{c_j(T)}{v-\eta_j(T)},\qquad N\le2M.
\end{equation}
The poles $\{\eta_j\}$ and residues $\{c_j\}$ are extracted by a
second Chebyshev pass on the Laurent coefficients
\begin{equation}\label{eq:bk}
b_k(T):=\bigl[v^{-(k+1)}\bigr]\,\Phi_M(v,T),\qquad k\ge0.
\end{equation}
Since $a_j(0)=a_j(-i)=0$ for all $j$, we have $d_k(0)=d_k(-i)=0$,
hence
\begin{equation}\label{eq:norm}
\Phi_M(0,T)=\Phi_M(-i,T)=0,\qquad
\varphi_M(0)=\varphi_M(-i)=1,
\end{equation}
where $\varphi_M:=\exp\Phi_M$. The identities
\begin{equation}
P(-\bar v)=\overline{P(v)},\qquad Q(-\bar v)=\overline{Q(v)},
\qquad R\in\R
\end{equation}
propagate to
\begin{equation}\label{eq:sym}
\Phi_M(-\bar v,T)=\overline{\Phi_M(v,T)},
\end{equation}
so the poles come in conjugate pairs, lie strictly off $\R$, and
\begin{equation}\label{eq:pstar}
p^*:=\min_j\bigl(-\Imag\,\eta_j\bigr)>0.
\end{equation}

% ------------------------------------------------------------------
\section{Main theorem}
% ------------------------------------------------------------------

\begin{theorem}[Exact rough Heston call price]\label{thm:main}
With the notation above, define
\begin{align}
H_j(v) &:= \exp\!\left(\sum_{i\ne j}\frac{c_i}{v-\eta_i}\right),
\label{eq:Hj}\\
G_j(v) &:= \frac{e^{-ivk}H_j(v)}{v(v+i)},
\label{eq:Gj}\\
\mathcal{R}_j &:= \sum_{n=0}^{\infty}
\frac{c_j^{n+1}\,G_j^{(n)}(\eta_j)}{(n+1)!\,n!}.
\label{eq:Rj}
\end{align}
For any $\alpha\in(1,p^*)$,
\begin{equation}\label{eq:price}
\boxed{\;
C_M=\begin{cases}
(S_0-Ke^{-rT})
-Ke^{-rT}\,\Real\!\left[i\displaystyle
\sum_{\Imag\,\eta_j>-\alpha}\!\!\mathcal{R}_j\right],
& k\le0,\\[1.8ex]
-\,Ke^{-rT}\,\Real\!\left[i\displaystyle
\sum_{\Imag\,\eta_j<-\alpha}\!\!\mathcal{R}_j\right],
& k>0.
\end{cases}\;}
\end{equation}
The series~\eqref{eq:Rj} converges absolutely. $C_M\to C$
at the geometric rate determined by the logarithmic capacity of
the singular set of $\widetilde\Phi(\cdot;v)$.
\end{theorem}

\begin{proof}
\textbf{Step 1 (Payoff transform).}
For $\Imag v<-1$, integrating over $x>k$ and using $S_0 e^{k+rT}=K$
gives
\begin{equation}
\int_{\R}e^{-ivx}(S_0 e^{x+rT}-K)^+\,dx
=-\frac{Ke^{-ivk}}{v(v+i)}.
\end{equation}
Risk-neutral pricing and Plancherel give
\begin{equation}\label{eq:parseval}
C_M=-\frac{Ke^{-rT}}{2\pi}\int_{\Imag v=-\alpha}
\frac{e^{-ivk}\varphi_M(v)}{v(v+i)}\,dv.
\end{equation}
The contour shift from $\Imag v=-1-\varepsilon$ to $\Imag v=-\alpha$
is valid since $\varphi_M$ is holomorphic and bounded on
$-p^*<\Imag v<0$ and
\begin{equation}
\Phi_M(v,T)=O(1/|v|)\quad\text{as }|v|\to\infty
\end{equation}
on horizontal lines, so Jordan's lemma applies.

\textbf{Step 2 (Contour closure).}
Set
\begin{equation}
F(v):=\frac{e^{-ivk}\varphi_M(v)}{v(v+i)}.
\end{equation}

\emph{Case $k\le0$:} since $|e^{-ivk}|=e^{(\Imag v)k}\to0$ as
$\Imag v\to+\infty$, close upward (counterclockwise).
Using~\eqref{eq:norm},
\begin{equation}
\Res_{v=0}F=-i,\qquad\Res_{v=-i}F=ie^{-k}.
\end{equation}
The kernel contribution is
\begin{equation}
-Ke^{-rT}\cdot i\bigl[(-i)+ie^{-k}\bigr]
=Ke^{-rT}(e^{-k}-1)=S_0-Ke^{-rT}.
\end{equation}
The model poles with $\Imag\,\eta_j>-\alpha$ contribute
$-iKe^{-rT}\sum\mathcal{R}_j$; taking $\Real$ gives the first
branch of~\eqref{eq:price}.

\emph{Case $k>0$:} since $|e^{-ivk}|\to0$ as $\Imag v\to-\infty$,
close downward (clockwise). Kernel poles $v=0,-i$ are above
$\Imag v=-\alpha$ and are not enclosed. Hence
\begin{equation}
\int_{\Imag v=-\alpha}F\,dv
=-2\pi i\sum_{\Imag\,\eta_j<-\alpha}\mathcal{R}_j,
\end{equation}
and therefore
\begin{equation}
C_M
=-\frac{Ke^{-rT}}{2\pi}\cdot(-2\pi i)
\sum_{\Imag\,\eta_j<-\alpha}\mathcal{R}_j
=-Ke^{-rT}\,\Real\!\left[i\sum_{\Imag\,\eta_j<-\alpha}\mathcal{R}_j\right],
\end{equation}
the second branch of~\eqref{eq:price}.

\textbf{Step 3 (Residue series).}
Near $\eta_j$ write $\varphi_M(v)=e^{c_j/(v-\eta_j)}H_j(v)$
with $H_j$ holomorphic and nonzero there. Expanding,
\begin{equation}
e^{c_j/(v-\eta_j)}=\sum_{p=0}^{\infty}\frac{c_j^p}{p!}(v-\eta_j)^{-p},
\qquad
G_j(v)=\sum_{n=0}^{\infty}\frac{G_j^{(n)}(\eta_j)}{n!}(v-\eta_j)^n.
\end{equation}
The coefficient of $(v-\eta_j)^{-1}$ requires $p=n+1$,
giving~\eqref{eq:Rj}. Setting $\rho:=\min_{i\ne j}|\eta_i-\eta_j|$,
Cauchy's inequality gives
\begin{equation}
\sum_{n=0}^{\infty}\frac{|c_j|^{n+1}|G_j^{(n)}(\eta_j)|}{(n+1)!\,n!}
\le M_\rho|c_j|\,e^{|c_j|/\rho}<\infty,
\end{equation}
so the series converges absolutely.

\textbf{Step 4 (Reality and convergence).}
Conjugate-pair symmetry~\eqref{eq:sym} gives
$\mathcal{R}_{\bar\eta_j}=\overline{\mathcal{R}_{\eta_j}}$, so
$C_M\in\R$. Stahl's theorem gives $\Phi_M\to\Phi$ uniformly on
$\Imag v=-\alpha$; Lipschitzness of $\exp$ then gives
$\varphi_M\to\varphi$ and $C_M\to C$ at the stated geometric rate.
\end{proof}

% ------------------------------------------------------------------
\section{Algorithm}
% ------------------------------------------------------------------

The following steps compute $C_M$ to prescribed precision.
No quadrature, FFT, or linear system solve is required.

\begin{enumerate}

\item \textbf{M\"untz--Tau recurrence.}
For $j=1,\ldots,2M+2$, compute $a_j(v)$
from~\eqref{eq:a1}--\eqref{eq:aj} as polynomials in $v$
with coefficients $P(v)$, $Q(v)$, $R$.

\item \textbf{Cumulant moments.}
For $k=0,\ldots,2M$, form
\begin{equation}
d_k(v)=u_k\,a_k(v)+w_k\,a_{k+1}(v)
\end{equation}
using the weights from Section~2.

\item \textbf{First Chebyshev pass.}
Evaluate $d_k(v)$ at required $v$-values.
Run the Chebyshev/Stieltjes algorithm on
$\{d_0(v),\ldots,d_{2M}(v)\}$ for $M+1$ steps to obtain
$\alpha_n(v)$, $\beta_n(v)$ for $n=0,\ldots,M$.

\item \textbf{FOPS and Pad\'e element.}
Evaluate the TTRR~\eqref{eq:TTRR} at $z=T^\mu$ to obtain
$p_{M+1}(T^\mu;v)$ and $p_{M+1}^{(1)}(T^\mu;v)$. Form
\begin{equation}
\Phi_M(v,T)=T\cdot d_0(v)\cdot
\frac{p_{M+1}^{(1)}(T^\mu;v)}{p_{M+1}(T^\mu;v)}.
\end{equation}
Check the bound~\eqref{eq:aposteriori}; increment $M$ if not satisfied.

\item \textbf{Laurent coefficients.}
Expand $\Phi_M(v,T)$ at $v=\infty$ to extract
\begin{equation}
b_k(T)=\bigl[v^{-(k+1)}\bigr]\,\Phi_M(v,T),\qquad k=0,\ldots,2M.
\end{equation}

\item \textbf{Second Chebyshev pass.}
Run the Chebyshev/Stieltjes algorithm on $\{b_k(T)\}$ to obtain
Jacobi matrix $J_N^{(v)}$; eigendecompose to get poles
$\{\eta_j(T)\}$ and residues $\{c_j(T)\}$.

\item \textbf{Residue sums and price.}
For each $j$, compute $G_j^{(n)}(\eta_j)$ for $n=0,1,2,\ldots$
until the tail bound
\begin{equation}
\frac{M_\rho\,|c_j|^{n+2}}{\rho^n\,(n+2)!}
\end{equation}
falls below the required precision. Assemble $\mathcal{R}_j$
from~\eqref{eq:Rj} and evaluate~\eqref{eq:price}.

\end{enumerate}

\end{document}
